%%%%%%%%%%%%%%%%%%%%%%%%%%%%%%%%%%%%%%%%%%%%%%%%%%%%%%%%%%%%%%%%%%%%%%%%%%%%%%%
% Nom           : 01_chap1.tex
% Rôle          : Rendu Chapitre 1 — Généralités sur les ordinateurs
% Auteur        : Maguendra Codandamourty (825041861)
% Cours         : Architecture des Machines et des Ordinateurs — DN22EM03
%%%%%%%%%%%%%%%%%%%%%%%%%%%%%%%%%%%%%%%%%%%%%%%%%%%%%%%%%%%%%%%%%%%%%%%%%%%%%%%

\chapter{Généralités sur les ordinateurs}

%%%%%%%%%%%%%%%%%%%%%%%%%%%%%%%%%%%%%%%%%%
% EXERCICE 1.2 — Au Choix
%%%%%%%%%%%%%%%%%%%%%%%%%%%%%%%%%%%%%%%%%%
\setcounter{section}{1} %permet de démarrer à 1.2
\section{Définitions du dictionnaire (Au Choix)}

\image{Énoncé --- Exercice 1.2 (Au Choix)}{1.1}{ex1_2_enonce.png}{fig:enonce12}

\begin{graybox}
    \textbf{Réponse :}
    Pour répondre à cette question, nous avons consulté le site Internet de
    l'académie Française. En effet, il s'agit d'une source fiable et bien souvent
    utilisé par les enseignants universitaires. J'ai choisi cette source aussi
    pour la qualité et la profondeur des définitions.

    \newpage

    \image{Site Internet de l'Académie Française}{1.2}{academie.png}{fig:academie}
    \subsection{Définitions (Dictionnaire de l'Académie française, 9\ieme{} éd.)}

    \subsubsection{Ordinateur}

    \textit{Nom masculin.} \parencite{acfr_ordinateur}

    \textbf{Étymologie :} \textit{xv}\textsuperscript{e} siècle, au sens de
    « celui qui institue quelque chose » ; \textit{xx}\textsuperscript{e}
    siècle, au sens actuel. Emprunté du latin \textit{ordinator}, « celui qui
    règle, met en ordre ; ordonnateur ».

    Équipement informatique comprenant les organes nécessaires à son
    fonctionnement autonome, qui assure, en exécutant les instructions d'un
    ensemble structuré de programmes, le traitement rapide de données codées
    sous forme numérique qui peuvent être conservées et transmises. La mise en
    réseau d'ordinateurs de bureau. Ordinateur portable. Ordinateur de poche.
    \textit{Microordinateur}, voir ce mot. Publication assistée par ordinateur.
    Les ordinateurs d'un centre de calcul (on dit aussi Calculateur numérique).

    \subsubsection{Programme}

    \textit{Nom masculin.} \parencite{acfr_programme}

    \textbf{Étymologie :} \textit{xvii}\textsuperscript{e} siècle. Emprunté
    du grec \textit{programma}, « ordre du jour, inscription », dérivé de
    \textit{prographein}, composé de \textit{pro}, « en avant, devant », et
    \textit{graphein}, « écrire ».

    \begin{enumerate}
        \item Feuillet, livret, affiche, etc. qui annonce la tenue d'un
        spectacle, d'une fête, d'une cérémonie, et donne le détail de leur
        déroulement ; ce qui est ainsi annoncé. Un programme affiché sur une
        colonne Morris. L'ouvreuse donne le programme. Ce morceau n'était pas
        au programme. Le programme d'une saison théâtrale, d'un festival.
        Loc. fam. \textit{Le programme des réjouissances}, ce qu'il est prévu
        de faire.
        \begin{itemize}
            \item Spécialement. Ensemble des émissions diffusées par une chaîne
            de télévision ou une station de radio ; par métonymie, feuillet,
            brochure en présentant la liste et les horaires. Le programme d'une
            radio, d'une télévision. La nouvelle grille des programmes.
            Consulter le programme de la semaine.
        \end{itemize}

        \item \textit{Marque de domaine : enseignement.} Ensemble de ce qui
        doit être enseigné à un élève d'une classe ou d'un cycle donnés ou de
        ce sur quoi porte un cours, un examen, un concours. Les programmes du
        primaire, du secondaire. La réforme des programmes scolaires. Une
        pièce de Molière figure au programme de l'agrégation de lettres cette
        année. Cette question est au programme, hors programme.

        \item Suite d'actions qu'on s'impose d'accomplir dans un but donné ;
        plan établi à l'avance. Un programme de travail. Définir un programme
        d'action. Tenir, remplir son programme, réaliser tout ce qu'on s'était
        promis de faire. Un programme d'investissement ambitieux.
        \begin{itemize}
            \item Loc. \textit{Vaste programme !} par allusion à une repartie
            plaisante du général de Gaulle. \textit{C'est tout un programme !}
            (souvent iron.), se dit d'une annonce qui étonne, qui laisse
            présager une suite surprenante.

            \item Spécialement. \textit{Marque de domaine : politique.}
            Ensemble des principes, des réformes, des mesures, etc. qu'un
            gouvernement, un parti ou un homme politique propose de mettre en
            application. Un programme économique, social. Le Premier ministre
            a exposé le programme du gouvernement. Contrat de programme,
            conclu entre l'État et des entreprises privées ou publiques, et
            qui permet d'harmoniser la politique du gouvernement et la gestion
            de ces entreprises.

            \item \textit{Marque de domaine : sports.} En patinage artistique,
            chacune des deux épreuves sur lesquelles les patineurs sont jugés
            en compétition. Le programme technique ou programme court, dont
            les figures sont imposées, s'oppose au programme libre ou programme
            long.

            \item \textit{Marque de domaine : architecture, bâtiment.}
            Document que le maître d'ouvrage fournit à un architecte et qui
            énumère ses exigences, ainsi que les caractéristiques du site
            choisi, les modalités de financement, etc. Le programme
            architectural d'un auditorium.

            \item Par extension. \textit{Marque de domaine : électronique,
            informatique.} Suite d'instructions rédigées dans un langage
            spécifique, qui commande à un dispositif, à un appareil ou à un
            système informatique d'exécuter une tâche donnée ; ensemble
            d'opérations, souvent cycliques, ainsi commandées. Concevoir,
            rédiger un programme. Le programme d'une machine-outil. Un
            programme d'exploitation, de contrôle. Programme de service, qui,
            dans un ordinateur, n'est pas destiné à traiter l'information mais
            à augmenter les possibilités initiales du système.

            \item \textit{Marque de domaine : musique.} Musique à programme,
            musique instrumentale ou vocale qui dépeint sous une forme sonore
            un élément non musical, le plus souvent un texte. La
            \textit{Symphonie fantastique}, d'Hector Berlioz, qui évoque un
            épisode de la vie de l'artiste, relève de la musique à programme.
        \end{itemize}
    \end{enumerate}

    \subsubsection{Donnée}

    \textit{Nom féminin.} \parencite{acfr_donnee}

    \textbf{Étymologie :} \textit{xiii}\textsuperscript{e} siècle, au sens de
    « distribution, aumône » ; \textit{xviii}\textsuperscript{e} siècle, comme
    terme de mathématiques. Participe passé féminin substantivé de
    \textit{donner} au sens d'« indiquer, dire ».

    \begin{enumerate}
        \item Fait ou principe indiscuté, ou considéré comme tel, sur lequel
        se fonde un raisonnement ; constatation servant de base à un examen,
        une recherche, une découverte. Les données de la science. Les données
        de l'expérience. Données statistiques. Ma théorie s'appuie sur des
        données précises. C'est un raisonnement fondé sur des données
        incertaines. Dans cette affaire, il est essentiel de connaître la
        donnée de départ. Par extension. Idée principale, thème constitutif
        qui est à l'origine d'une œuvre littéraire. La donnée d'une tragédie.

        \item \textit{Marque de domaine : psychologie.} Ce qui est connu
        immédiatement par le sujet, indépendamment de toute élaboration de
        l'esprit, par opposition à ce qui est connu par induction ou déduction,
        par raisonnement, par calcul. \textit{Titre célèbre :} Essai sur les
        données immédiates de la conscience, d'Henri Bergson (1889).

        \item \textit{Marque de domaine : mathématiques.} Souvent au pluriel.
        Chacune des quantités ou propriétés mentionnées dans l'énoncé d'un
        problème et qui permettent de le résoudre.

        \item \textit{Marque de domaine : informatique.} Représentation d'une
        information sous une forme conventionnelle adaptée à son exploitation.
        \textit{Le traitement automatique des données. Une banque, une base de
        données.}
    \end{enumerate}

    \subsubsection{Information}

    \textit{Nom féminin.} \parencite{acfr_information}

    \textbf{Étymologie :} \textit{xiii}\textsuperscript{e} siècle, au sens
    d'« enquête faite en matière criminelle ». Dérivé d'\textit{informer}.

    \begin{enumerate}
        \item Action d'informer ou de s'informer. Mission d'information.
        Voyage d'information. Réunion d'information. Un service d'information.
        Une note distribuée pour information.
        \begin{itemize}
            \item Par métonymie. Renseignement qu'on donne ou qu'on obtient.
            La recherche de l'information. Une information de source autorisée.
            Une information digne de foi. Communiquer une information. Souvent
            au pluriel. Des informations confidentielles. Aller aux
            informations, recueillir des informations. Prendre des informations
            sur quelqu'un. Manquer d'informations. Vérifier ses informations.

            \item Au singulier, dans un sens collectif. Ensemble de données, de
            connaissances réunies sur un sujet. Il s'est lancé dans cette
            entreprise avec une information insuffisante. Ces travaux témoignent
            d'une information très étendue.

            \item \textit{Marque de domaine : droit.} Enquête menée par le
            juge d'instruction ou par des officiers de police judiciaire en
            vue d'établir la preuve d'une infraction et d'en découvrir les
            auteurs. Information judiciaire. Procéder à une information.
            Demander un supplément d'information. Ouvrir, clore une
            information. Information contre X.
        \end{itemize}

        \item Spécialement. Action de porter des nouvelles à la connaissance
        du public, de faire part des évènements, des faits marquants de
        l'actualité. Un grand quotidien d'information. Liberté d'information.
        Les métiers de l'information. Les techniques, les moyens d'information.
        \begin{itemize}
            \item Par métonymie. Fait ou évènement dont font état une agence
            ou un organe de presse, la radio, la télévision. Informations
            générales, politiques, économiques, sportives. Une information de
            dernière heure. Bulletin d'information. Par extension. Au pluriel.
            Informations radiodiffusées, télévisées ou, simplement,
            \textit{informations}, émission de radio ou de télévision donnant
            les nouvelles du jour.
        \end{itemize}

        \item \textit{Marque de domaine : informatique.} Élément de
        connaissance traduit par un ensemble de signaux selon un code
        déterminé, en vue d'être conservé, traité ou communiqué.
        \textit{Traitement de l'information}, emploi d'ordinateurs en vue
        d'effectuer des opérations logiques et mathématiques complexes à des
        fins scientifiques, administratives, etc. Sciences de l'information,
        disciplines concernant l'utilisation de ces techniques dans divers
        domaines professionnels.
        \begin{itemize}
            \item Par analogie. \textit{Marque de domaine : biologie.}
            Information génétique, ensemble des messages qui sont portés par
            les gènes et fournissent à un organisme vivant les règles de sa
            structure, de son développement et de son action. Les molécules
            d'acide désoxyribonucléique constituent le support de
            l'information génétique.
        \end{itemize}
    \end{enumerate}

    \subsubsection{Connaissance}

    \textit{Nom féminin.} \parencite{acfr_connaissance}

    \textbf{Étymologie :} \textit{xii}\textsuperscript{e} siècle,
    \textit{conoisance}, « acte de connaître, idée, notion de quelque chose ».
    Dérivé du radical du participe présent de \textit{connaître}.

    \textbf{I.} Exercice de la faculté par laquelle on connaît et distingue
    les objets, ainsi que les actes ou états du sujet.

    \begin{enumerate}
        \item Acte de l'esprit par lequel on se représente, définit ou
        comprend un objet. Les voies de la connaissance. Connaissance
        intuitive, rationnelle, empirique, sensible. Une connaissance claire et
        distincte. La connaissance scientifique. Théorie de la connaissance.
        Critique de la connaissance.

        \item Le fait d'être informé ou de s'informer, d'apprendre quelque
        chose. Cela est venu à la connaissance des autorités. Porter un
        règlement à la connaissance du public. J'ai eu, je n'ai pas eu
        connaissance de cet évènement, de cette affaire.
        \begin{itemize}
            \item Expr. \textit{À ma connaissance}, autant que je sache.
            \textit{À ma connaissance, il est déjà de retour.}
        \end{itemize}

        \item Idée, notion que l'on a d'une personne ou d'une chose,
        représentation que l'on s'en fait. La connaissance de Dieu. La
        connaissance du bien et du mal. La connaissance des hommes, du cœur
        humain. La connaissance de l'avenir.

        \item Sentiment de sa propre existence ; plein exercice de ses
        facultés. Perdre connaissance, tomber sans connaissance, s'évanouir.
        Elle est restée longtemps sans connaissance. Il a repris rapidement
        connaissance. Il a conservé toute sa connaissance.

        \item Ce que l'on connaît par l'étude, l'expérience ou par tout autre
        moyen d'information. Avoir une connaissance théorique, pratique. Avoir
        d'un sujet une connaissance approximative, vague, fragmentaire,
        précise, exhaustive, exacte. Il n'a qu'une connaissance superficielle
        du dossier. Il a une grande connaissance des langues, une solide
        connaissance de la musique.
        \begin{itemize}
            \item Au pluriel. Ensemble de ce que l'on a appris ; savoir,
            acquis. Avoir de grandes, de profondes, de vastes connaissances.
            Des connaissances sommaires, élémentaires, fragmentaires. Acquérir,
            amasser, accumuler les connaissances. Faire montre de ses
            connaissances. \textit{Marque de domaine : enseignement.} Contrôle
            continu des connaissances, mode d'évaluation des acquisitions de
            l'élève par des contrôles partiels et fréquents. \textit{Titre
            célèbre :} Connaissance de l'Est, de Paul Claudel (1900 et 1907).
        \end{itemize}

        \item Spécialement.
        \begin{itemize}
            \item \textit{Marque de domaine : droit.} Le droit de connaître
            de certaines affaires, le droit de juger. Attribuer à un tribunal
            la connaissance de certaines causes. Expr. \textit{En connaissance
            de cause}, voir Cause.

            \item \textit{Marque de domaine : marine.} Avoir connaissance d'une
            côte, d'une île, l'apercevoir et l'identifier.

            \item \textit{Marque de domaine : vènerie.} Avoir connaissance
            d'une bête, en apercevoir les traces. Au pluriel. Traces laissées
            par le pied de l'animal, donnant des indications sur son âge, sa
            taille, etc.

            \item \textit{Marque de domaine : astronomie.} \textit{Connaissance
            des temps}, volume de tables, publié chaque année par le Bureau des
            longitudes, et donnant à l'avance les éléments variables des
            différents astres.
        \end{itemize}
    \end{enumerate}

    \medskip
    \textbf{II.} Relation, rapport qu'on établit ou qu'on entretient avec une
    ou plusieurs autres personnes.

    \begin{enumerate}
        \item Le fait d'entrer en rapport. Faire connaissance, se lier, entrer
        en relation. Il a fait connaissance avec ses voisins, il a fait la
        connaissance de ses voisins. Vous devriez faire sa connaissance. Nous
        avons renoué, refait connaissance après une longue séparation. Elle ne
        lie pas facilement connaissance. Fig. Faire connaissance avec la faim,
        la misère, la solitude, la prison, en faire une première expérience.
        \begin{itemize}
            \item Loc. \textit{De connaissance}, que l'on connaît. Je vois
            là-bas une figure de connaissance. Nous sommes entre gens de
            connaissance, qui se connaissent. Expr. Être, se trouver en pays
            de connaissance, dans des lieux familiers, parmi des gens que l'on
            fréquente habituellement, dans un milieu où l'on s'intéresse aux
            mêmes sujets.
        \end{itemize}

        \item Par métonymie. Fam. Personne que l'on connaît, avec qui l'on est
        en relation. Ce n'est pas un ami, mais une simple connaissance, une
        vague connaissance. Nous sommes de vieilles connaissances.
    \end{enumerate}

    \subsection{Conclusion}

    Il est intéressant de définir des termes que l'on utilise tous les jours et
    notamment de voir l'évolution de ces termes. J'apprécie particulièrement les
    définitions du dictionnaire de l'académie française, car elles sont, à mon sens,
    très complètes et elles permettent de recenser toutes les acceptions d'un terme.
    Les descriptions étymologiques sont très parlantes et nous renvoient aux significations
    originelles des termes.

    \subsection{Références}
    \begin{itemize}
        \item[1] Dictionnaire de l'académie Française - Ordinateur \\
         \url{https://www.dictionnaire-academie.fr/article/A9O0665},\\
         Consulté le 20/05/2026
        \item[2] Dictionnaire de l'académie Française - Programme \\
        \url{https://www.dictionnaire-academie.fr/article/A9P4504},\\
         Consulté le 20/05/2026
        \item[3] Dictionnaire de l'académie Française - Donnée\\
        \url{https://www.dictionnaire-academie.fr/article/A9D3040},\\
         Consulté le 20/05/2026
        \item[4] Dictionnaire de l'académie Française - Information\\
        \url{https://www.dictionnaire-academie.fr/article/A9I1218},\\
         Consulté le 20/05/2026
        \item[5] Dictionnaire de l'académie Française - Connaissance\\
        \url{https://www.dictionnaire-academie.fr/article/A9C3633},\\
         Consulté le 20/05/2026,
    \end{itemize}

\end{graybox}

\newpage

%%%%%%%%%%%%%%%%%%%%%%%%%%%%%%%%%%%%%%%%%%
% EXERCICE 1.3 — À Rendre
%%%%%%%%%%%%%%%%%%%%%%%%%%%%%%%%%%%%%%%%%%

\section{Instruction codée par 4, 4, 4, 4 (À Rendre)}

\image{Énoncé --- Exercice 1.3 (À Rendre)}{1.1}{ex1_3_enonce.png}{fig:enonce13}

\begin{graybox}
    \textbf{Réponse :}

    \subsection{Analyse de l'instruction}

    Selon le tableau jeu d'instructions du cours (tableau 1.3.4), le code opération 4 signifie 
    écrire à adr3 la somme des contenus de adr1 et adr2. 
    L’instruction 4, 4, 4, 4 signifie donc écrire la somme des contenus de l’adr4 et adr4
    puis stocker le résultat à l’adr4. L’adr4 contient la valeur 251. 
    Donc on aura 251 + 251 = 502. Le programme va écraser la valeur située 
    à l’adresse 4 (251) et la remplacer par 502. 
    

    \subsection{Conclusion}

    Cet exercice nous apprend comment additionner deux fois la même valeur sans 
    stocker une nouvelle fois la valeur à une autre adresse. On voit aussi
    que la modification du contenu à l'adresse 4 aura des répercussions sur les 
    futures instructions. Il faut donc être vigilant et bien comprendre la suite
    d'instructions que l'on exécute. 

\end{graybox}

\newpage

%%%%%%%%%%%%%%%%%%%%%%%%%%%%%%%%%%%%%%%%%%
% EXERCICE 1.4 — Au Choix
%%%%%%%%%%%%%%%%%%%%%%%%%%%%%%%%%%%%%%%%%%

\setcounter{section}{3}
\section{Instruction codée par 2, 4, 4, 4 et division par zéro (Au Choix)}

\image{Énoncé --- Exercice 1.4 (Au Choix)}{1.1}{ex1_4_enonce.png}{fig:enonce14}

\begin{graybox}
    \textbf{Réponse :}

    \subsection{Analyse de l'instruction 2, 4, 4, 4}
    Le code instruction 2 nous dit “ écrire à adr3 le résultat de la division des 
    contenus adr1 et adr2”. L’adr1 et adr2 continent la valeur 502 (suite à l’exercice 1.3).

    On aura donc $502 \div 502 = 1$. 
    
    L’adr4 contiendra donc la valeur 1.

    Dans le cas du remplacement de l’adr6 par 0, on a le codage suivant :
    2, adr5, adr6, adr7. On remplace donc le dénominateur par 0 (remplace l’adr6 par 0) 
   cela revient à diviser par 0, ce qui peut provoquer une erreur ou un arrêt du programme. 
   En effet, la division par 0 est mathématiquement impossible. Si le programme a anticipé 
   ce type d'erreur, alors il affichera un message d'erreur tel que "impossible de diviser par 0". 


    \subsection{Conclusion}

    Cet exercice pointe l'importance de vérifier la validité des opérandes avant de lancer
    le programme. Ce type d'erreur peut être fréquent dans les programmes si l'on ne met pas
    en place de garde-fou pour éviter cette situation. Il m'est arrivé bien souvent des erreurs
    de ce type, car j'ai du mal à suivre dans ma tête la valeur que vont prendre mes variables. 
    Je pense que c'est une compétence que je dois développer. Ce qui m'aide souvent sont les 
    exécuteur de programme "pas à pas" qui me permettent de vérifier que les variables prennent
    les valeurs attendues. Je commence à ajouter dans mon code, notamment en Python, des fonctions
    qui me permettent de mieux gérer les erreurs. 

\end{graybox}

\newpage

%%%%%%%%%%%%%%%%%%%%%%%%%%%%%%%%%%%%%%%%%%
% EXERCICE 1.6 — À Rendre
%%%%%%%%%%%%%%%%%%%%%%%%%%%%%%%%%%%%%%%%%%

\setcounter{section}{5}
\section{Instructions en boucle (À Rendre)}

\image{Énoncé --- Exercice 1.6 (À Rendre)}{1.1}{ex1_6_enonce.png}{fig:enonce16}

\begin{graybox}
    \textbf{Réponse :}

    \subsection{Décodage des instructions}

    La mémoire contient les valeurs suivantes à partir de l'adresse 2000 :

    \begin{center}
        \begin{tabular}{lll}
            \hline
            \textbf{Adresse} & \textbf{Contenu} & \textbf{Signification} \\
            \hline
            2000 & 3    & Code opération : continuer avec une autre instruction \\
            2001 & 2000 & La prochaine instruction se trouve à l'adresse 2000 \\
            2002 & 2001 & Ignoré (non utilisé par l'opération 3) \\
            2003 & 2002 & Ignoré (non utilisé par l'opération 3) \\
            \hline
        \end{tabular}
    \end{center}
    
    L’opération 3 demande de continuer avec une autre instruction. L’adresse suivante 2001 nous 
    renvoie vers l’adresse 2000 qui dont nous demande de continuer avec la prochaine instruction. 
    L’adresse 2002 et 2003 peuvent être ignoré si l’on compare au tableau 1.3.5 : 

    \begin{itemize}
        \item adresse 3017 : \code{3} → continuer avec une autre instruction
        \item adresse 3018 : \code{3005} → la prochaine instruction est à l'adresse 3005
        \item adresses 3019 et 3020 : ignorées
    \end{itemize}

    Le programme lit l'instruction à 2000 (code \code{3}), consulte le contenu
    de l'adresse 2001 (valeur 2000) et reprend l'exécution à l'adresse 2000.    
    Je pense que cela va nous donner une boucle infinie car l’instruction est constamment 
    renvoyée vers l’adresse 2000 qui demande de continuer avec une autre instruction. 


    \subsection{Conclusion}

    On voit dans cet exercice un autre type d'erreur qui peut se produire très facilement,
    à savoir la création d'une boucle infinie non intentionnel. C'est un type d'erreur qui  
    m'arrive avec notamment les boucle 'while'. Il est important de vraiment comprendre
    en profondeur la liste des instructions avant de lancer le code. J'ai vu dans un cours
    en ligne (CS50 Introduction to Computer Science \footnote{https://pll.harvard.edu/course/cs50-introduction-computer-science consulté le 20/05/26})
    qu'il était recommandé de parler à un petit canard en plastique afin de vraiment s'assurer
    que le code est "logique". J'ai trouvé l'idée amusante, et je pense mon temps à parler à 
    mon chien (qui semble bien intéressé par la programmation !). 

\end{graybox}

\newpage

%%%%%%%%%%%%%%%%%%%%%%%%%%%%%%%%%%%%%%%%%%
% EXERCICE 1.7 — À Rendre
%%%%%%%%%%%%%%%%%%%%%%%%%%%%%%%%%%%%%%%%%%

\section{Suite de Fibonacci (À Rendre)}

\image{Énoncé --- Exercice 1.7 (À Rendre)}{1.1}{ex1_7_enonce.png}{fig:enonce17}

\begin{graybox}
    \textbf{Réponse :}

    \subsection{Introduction}

    Pour répondre à cet exercice, je construis un tableau pour remplir les
    valeurs successives.

    \subsection{Tableau détaillé des itérations}

    \resizebox{\linewidth}{!}{%
    \begin{tabular}{|c|c|l|c|c|c|c|c|c|}
        \hline
        \textbf{Adresse} & \textbf{Contenu} & \textbf{Signification} & \textbf{it2} & \textbf{it3} & \textbf{it4} & \textbf{it5} & \textbf{it6} & \textbf{it7} \\
        \hline
        1997 & 1    & Placer                                                          & & & & & & \\
        \hline
        1998 & 1    & le nombre 1                                                     & & & & & & \\
        \hline
        1999 & 2999 & à l'adresse 2999                                                & & & & & & \\
        \hline
        2000 & 3000 & 0                                                               & & & & & & \\
        \hline
        2001 & 1    & Placer                                                          & & & & & & \\
        \hline
        2002 & 0    & le nombre 0                                                     & & & & & & \\
        \hline
        2003 & 3000 & à l'adresse 3000                                                & & & & & & \\
        \hline
        2004 & 2004 & 0                                                               & & & & & & \\
        \hline
        2005 & 1    & Placer                                                          & & & & & & \\
        \hline
        2006 & 1    & le nombre 1                                                     & & & & & & \\
        \hline
        2007 & 3001 & à l'adresse 3001                                                & & & & & & \\
        \hline
        2008 & 3001 & 0                                                               & & & & & & \\
        \hline
        2009 & 4    & Ajouter                                                         & & & & & & \\
        \hline
        2010 & 3000 & le contenu de l'adresse 3000 (donc 0)                          & 1       & 1       & 2       & 3       & 5       & 8        \\
        \hline
        2011 & 3001 & au contenu de l'adresse 3001 (donc 1)                          & 1       & 2       & 3       & 5       & 8       & 13       \\
        \hline
        2012 & 3002 & et placer le résultat à l'adresse 3002 : donc 0+1=1            & 1+1=2   & 1+2=3   & 2+3=5   & 3+5=8   & 5+8=13  & 8+13=21  \\
        \hline
        2013 & 2    & Diviser                                                         & & & & & & \\
        \hline
        2014 & 3001 & le contenu de l'adresse 3001 : (donc 1)                        & 1       & 2       & 3       & 5       & 8       & 13       \\
        \hline
        2015 & 2999 & par le contenu de l'adresse 2999 (donc 1)                      & 1       & 1       & 1       & 1       & 1       & 1        \\
        \hline
        2016 & 3000 & et placer le résultat à l'adresse 3000 (donc 1)                & 1       & 2/1=2   & 3/1=3   & 5/1=5   & 8/1=8   & 13/1=13  \\
        \hline
        2017 & 2    & Diviser                                                         & & & & & & \\
        \hline
        2018 & 3002 & le contenu de l'adresse 3002 (donc 1)                          & 2       & 3       & 5       & 8       & 13      & 21       \\
        \hline
        2019 & 2999 & par le contenu de l'adresse 2999 (1)                           & 1       & 1       & 1       & 1       & 1       & 1        \\
        \hline
        2020 & 3001 & et placer le résultat à l'adresse 3001 (donc 1)                & 2       & 3/1=3   & 5/1=5   & 8/1=8   & 13/1=13 & 21/1=21  \\
        \hline
        2021 & 3    & Continuer avec une autre instruction                            & & & & & & \\
        \hline
        2022 & 2009 & la prochaine instruction est à l'adresse 2009                  & & & & & & \\
        \hline
    \end{tabular}%
    }

    \subsection{Valeurs successives}

    Les valeurs successives des adresses sont :

    \begin{center}
        \begin{tabular}{ll}
            \hline
            \textbf{Adresse} & \textbf{Valeurs successives} \\
            \hline
            2999 & 1 \\
            3000 & 0, 1, 1, 2, 3, 5, 8 \\
            3001 & 1, 1, 2, 3, 5, 8, 13 \\
            3002 & 1, 2, 3, 5, 8, 13, 21 \\
            \hline
        \end{tabular}
    \end{center}

    Par itération successive, on obtient donc :
    $0, 1, 1, 2, 3, 5, 8, 13, \ldots$ à l’adresse 3000.
    Il s’agit de la \textbf{suite de Fibonacci}.



    \subsection{Conclusion}
    Cet exercice nous renvoie à la définition même de ce que signifie le mot
    ordinateur : ordonner un certain nombre de tâche. Il nous permet de comprendre
    la logique de l'ordinateur et le changement de valeur successive. Le tableau
    que j'ai conçu m'a beaucoup aidé à comprendre la succession des tâches. Je 
    trouve que c'est intéressant que décortiquer un ordre donné étape par étape, cela 
    peut aider notamment lors du débugage. C'est intéressant de voir que l'on peut
    créer des suites de nombre bien connus avec une succession d'ordres. 
    On apprend dans un article Wikipédia\footnote{\url{https://fr.wikipedia.org/wiki/Suite_de_Fibonacci}} que la suite de Fibonacci est liée au fameux nombre d'or.

    

\end{graybox}
